%% ma: L12-B08-0130
%% lop: 12
%% chuong: 2
%% bai: 8
%% dang: 12.8.46
%% loai: TN4
%% muc: VD
%% dap_an: "C"
%% nguon: "CĐ2.3.pdf giáo viên gửi 27/10/2026 (M07, Lê Minh Khôi – Nguyễn Anh Khôi, Chương II Chủ đề 3 Biểu thức tọa độ của các phép toán vectơ), Câu 45"
%% kien_thuc: "Bài 8 (khoảng cách hai điểm)"
%% trang_thai: chinh-thuc
%% ghi_chu: "Giáo viên đã duyệt 27/10/2026"
\begin{ex}
Trong không gian $Oxyz$, cho hình bình hành $MNPQ$ có $M(-4;3;3)$, $N(4;-4;2)$ và $P(3;6;-1)$. Biết chu vi của hình bình hành $MNPQ$ bằng $\sqrt a+\sqrt b$ với $a,b\in\mathbb N^*$ và $a>b$. Tính giá trị của biểu thức $S=a-b$.
\choice{$S=4$}{$S=8$}{\True $S=16$}{$S=32$}
\loigiai{$\overrightarrow{MN}=(8;-7;-1)$, $MN=\sqrt{114}$; $\overrightarrow{NP}=(-1;10;-3)$, $NP=\sqrt{110}$.

Chu vi $2\sqrt{114}+2\sqrt{110}=\sqrt{456}+\sqrt{440}$, $S=16$.}
\end{ex}
